\subtituloA{Limitaciones y seguridad del entorno cliente}

\definicion{Sandbox}{Entorno de ejecución aislado en el navegador que impide a JavaScript leer o escribir archivos locales, lanzar programas o modificar la configuración del navegador.}

\lista{
  \textbf{Restricción de archivos:} No puede acceder al disco duro del usuario. |
  \textbf{Política del mismo origen (SOP):} Un script de dominio A no puede acceder a recursos de dominio B. El origen se define como protocolo + dominio + puerto. |
  \textbf{CORS:} Error que ocurre al acceder a APIs ajenas debido a la aplicación de la política SOP; requiere permiso explícito del servidor mediante cabeceras. |
  \textbf{Vulnerabilidad XSS:} Riesgo si se inserta texto no limpiado de un usuario en la página, permitiendo la ejecución de scripts maliciosos.
}

\nota{Para el examen}{
Es falso que HTML gestione páginas dinámicas; la dinámica requiere lenguajes de script en el servidor y JavaScript en el cliente. JavaScript no puede escribir archivos en el servidor debido al modelo sandbox.
}

\begin{codigo}
 // Ejemplo de manipulación dinámica de estilo en cliente 
const elemento = document.getElementById('titulo'); 
elemento.style.color = 'blue';
\end{codigo}


